% !TEX program = xelatex
\documentclass[aspectratio=169,11pt,handout]{beamer}
\usetheme{Madrid}
\usepackage{fontspec}
\setsansfont{DejaVu Sans}
\setmainfont{DejaVu Serif}
\usepackage{polyglossia}
\setdefaultlanguage{turkish}
\usepackage{amsmath,amssymb}
\usepackage{booktabs}
\usepackage{multicol}
\usepackage{xcolor}
\usepackage{tikz}

% ---------- Renkler ----------
\definecolor{anamavi}{RGB}{0,45,98}
\definecolor{cozumyesil}{RGB}{0,110,50}
\definecolor{uyarikirmizi}{RGB}{180,20,20}
\setbeamercolor{structure}{fg=anamavi}
\setbeamercolor{frametitle}{bg=anamavi,fg=white}
\setbeamercolor{title}{bg=anamavi,fg=white}
\setbeamercolor{palette primary}{bg=anamavi,fg=white}
\setbeamercolor{palette secondary}{bg=anamavi!85!black,fg=white}
\setbeamercolor{palette tertiary}{bg=anamavi!70!black,fg=white}
\setbeamercolor{block title}{bg=anamavi,fg=white}
\setbeamercolor{block body}{bg=anamavi!7}
\setbeamercolor{block title example}{bg=cozumyesil,fg=white}
\setbeamercolor{block body example}{bg=cozumyesil!8}
\setbeamercolor{block title alerted}{bg=uyarikirmizi,fg=white}
\setbeamercolor{block body alerted}{bg=uyarikirmizi!7}
\setbeamertemplate{navigation symbols}{}
\setbeamertemplate{blocks}[rounded][shadow=false]

% ---------- Soru / Çözüm kalıbı ----------
% \soru{başlık}{soru metni}{çözüm metni}: önce soru, bir tıklamayla çözüm
\newcommand{\soru}[3]{%
\begin{frame}{#1}
  \begin{block}{Soru}\large #2\end{block}
  \pause
  \begin{exampleblock}{Çözüm}\large #3\end{exampleblock}
\end{frame}}
\newcommand{\ab}[1]{\overline{#1}}

\title[MMAT105 -- Hafta 4]{MMAT105 Genel Matematik}
\subtitle{Hafta 4: Rasyonel Sayılar, Ondalık Sayılar, Dört İşlem ve Sıralama}
\author{}
\date{2026--2027 Güz Dönemi}

\begin{document}

% =====================================================
\begin{frame}
\titlepage
\end{frame}

% (iletişim çerçevesi kaldırıldı)

\begin{frame}{Geçen Hafta Hatırlatma}
\begin{columns}[t]
\column{0.6\textwidth}
\textbf{3. hafta: bölünebilme, EBOB, EKOK}
\begin{enumerate}
\item Üç basamaklı \(\ab{37a}\) sayısı hem \(3\) hem \(4\) ile tam bölünebildiğine göre \(a\) rakamı kaçtır?
\item \(\text{EBOB}(24,36)\) ve \(\text{EKOK}(24,36)\) kaçtır?
\item \(12\) ve \(18\) ile bölündüğünde her ikisinde de \(5\) kalanını veren iki basamaklı en küçük doğal sayı kaçtır?
\end{enumerate}
\column{0.36\textwidth}
\pause
\textbf{Cevaplar:}
\begin{enumerate}
\item \(a=2\) \\ {\small (\(4\): \(a\in\{2,6\}\); \ \(3\): \(a\in\{2,5,8\}\))}
\item \(12\) \ ve \ \(72\)
\item \(\text{EKOK}+5=36+5=41\) \\ {\small (\(5\) de sağlar ama tek basamaklı)}
\end{enumerate}
\end{columns}
\vspace{4pt}
\small Bu hafta EKOK'u \textbf{payda eşitlemede}, EBOB'u \textbf{sadeleştirmede} kullanacağız.
\end{frame}

\begin{frame}{Bugün}
\begin{columns}[t]
\column{0.32\textwidth}
\begin{block}{Ders 1}
\begin{enumerate}
\item Rasyonel sayı, kesir türleri
\item Genişletme, sadeleştirme
\item Sayı doğrusu
\item Kesirlerde dört işlem
\item Merdiven kesirler
\end{enumerate}
\end{block}
\column{0.32\textwidth}
\begin{block}{Ders 2}
\begin{enumerate}\setcounter{enumi}{5}
\item Ondalık \(\leftrightarrow\) kesir
\item Ondalıklarda dört işlem
\item Devirli ondalık sayılar
\item Bilimsel gösterim
\end{enumerate}
\end{block}
\column{0.32\textwidth}
\begin{block}{Ders 3}
\begin{enumerate}\setcounter{enumi}{9}
\item Sıralama yöntemleri
\item Negatif rasyoneller
\item İki sayı arasında sayı
\item Karma işlem, kesir problemleri
\end{enumerate}
\end{block}
\end{columns}
\end{frame}

% =====================================================
\section{Ders 1: Rasyonel Sayılar ve Kesirlerde Dört İşlem}
\begin{frame}
\centering\Large\color{anamavi}\textbf{Ders 1}\\[6pt]\normalsize Rasyonel Sayılar, Kesir Türleri ve Kesirlerde Dört İşlem
\end{frame}

\subsection{Rasyonel sayılar}
\begin{frame}{Rasyonel Sayılar ve Kesir Türleri}
\begin{block}{Tanım}
\small
\[\mathbb{Q}=\left\{\tfrac{a}{b}\ :\ a,b\in\mathbb{Z},\ b\neq 0\right\},\qquad \mathbb{N}\subset\mathbb{Z}\subset\mathbb{Q}\qquad(\text{her } n\in\mathbb{Z} \text{ için } n=\tfrac{n}{1})\]
\vspace{-10pt}
\begin{itemize}\setlength{\itemsep}{1pt}
\item \(a,b>0\) iken \textbf{basit kesir}: \(a<b\); \quad \textbf{bileşik kesir}: \(a\ge b\) \quad (ör. \(\frac{3}{5}\) basit; \(\frac{7}{4},\ \frac{5}{5}\) bileşik)
\item \textbf{Tam sayılı kesir}: \(\displaystyle n\tfrac{p}{q}=n+\frac{p}{q}=\frac{n\cdot q+p}{q}\) \quad (ör. \(2\frac{1}{3}=\frac{7}{3}\))
\item \textbf{Genişletme / sadeleştirme}: \(\frac{a}{b}=\frac{a\cdot k}{b\cdot k}\ (k\neq 0)\); \ en sade biçim için pay ve payda EBOB'a bölünür.
\item \textbf{Denk kesirler}: \(\frac{a}{b}=\frac{c}{d}\iff a\cdot d=b\cdot c\) \quad (\(b,d\neq 0\)); \ ör. \(\frac{3}{4}=\frac{9}{12}=\frac{15}{20}\)
\end{itemize}
\end{block}
\begin{alertblock}{Sık yapılan hatalar}
\small\(\frac{0}{5}=0\) ama \(\frac{5}{0}\) \textbf{tanımsız}. \quad Toplamlarda sadeleştirme yapılmaz: \(\frac{2+3}{2+5}=\frac{5}{7}\neq\frac{3}{5}\). \quad \(-\frac{a}{b}=\frac{-a}{b}=\frac{a}{-b}\).
\end{alertblock}
\end{frame}

\soru{Rasyonel Sayılar -- Örnek 1}
{\(\dfrac{0}{7},\quad \dfrac{7}{0},\quad -\dfrac{3}{4},\quad 5,\quad 2{,}5,\quad \dfrac{12}{-4}\) \\[4pt] ifadelerinden hangileri rasyonel sayıdır? Hangileri tam sayıdır?}
{\(\frac{0}{7}=0\), \ \(5=\frac{5}{1}\), \ \(2{,}5=\frac{5}{2}\), \ \(-\frac{3}{4}\), \ \(\frac{12}{-4}=-3\) \quad \(\Rightarrow\) bunlar rasyoneldir.\\
\(\frac{7}{0}\) \textbf{tanımsızdır}, rasyonel sayı değildir.\\[4pt]
Tam sayı olanlar: \(\frac{0}{7}=0\), \ \(5\), \ \(\frac{12}{-4}=-3\).}

\soru{Kesir Türleri -- Örnek 2}
{\(\dfrac{3}{7},\quad \dfrac{9}{4},\quad \dfrac{6}{6},\quad \dfrac{11}{12},\quad \dfrac{13}{5}\) \\[4pt] kesirlerinden hangileri basit, hangileri bileşik kesirdir?}
{Basit (pay \(<\) payda): \(\dfrac{3}{7},\ \dfrac{11}{12}\) \quad (değeri \(1\)'den küçük)\\[4pt]
Bileşik (pay \(\ge\) payda): \(\dfrac{9}{4},\ \dfrac{6}{6},\ \dfrac{13}{5}\) \quad (değeri \(1\)'e eşit veya büyük)}

\soru{Tam Sayılı Kesir -- Örnek 3}
{\(3\dfrac{2}{5}\) sayısını bileşik kesir olarak, \(\dfrac{23}{4}\) kesrini tam sayılı kesir olarak yazınız.}
{\[3\frac{2}{5}=\frac{3\cdot 5+2}{5}=\frac{17}{5}\]
\(23=4\cdot 5+3\) \ (bölüm \(5\), kalan \(3\)):
\[\frac{23}{4}=5\frac{3}{4}\]}

\soru{Sadeleştirme -- Genişletme: Örnek 4}
{\(\dfrac{84}{126}\) kesrini en sade biçimde yazınız. \(\dfrac{3}{8}\) kesrini paydası \(56\) olacak şekilde genişletiniz.}
{\(\text{EBOB}(84,126)=42\):
\[\frac{84}{126}=\frac{84\div 42}{126\div 42}=\frac{2}{3}\]
\(56=8\cdot 7\) \(\Rightarrow\) \(7\) ile genişlet: \quad \(\dfrac{3}{8}=\dfrac{3\cdot 7}{8\cdot 7}=\dfrac{21}{56}\)}

\soru{Sayı Doğrusu -- Örnek 5}
{\(-\dfrac{7}{3}\) ve \(\dfrac{11}{4}\) sayıları hangi iki ardışık tam sayı arasındadır? Hangisi \(0\)'a daha yakındır?}
{\(-\frac{7}{3}=-2\frac{1}{3}\): \(-3\) ile \(-2\) arasında; \quad \(\frac{11}{4}=2\frac{3}{4}\): \(2\) ile \(3\) arasında.\\[2pt]
\begin{center}
\begin{tikzpicture}[scale=0.95]
\draw[->] (-3.6,0)--(3.6,0);
\foreach \x in {-3,...,3} {\draw (\x,0.1)--(\x,-0.1) node[below] {\small \(\x\)};}
\fill[uyarikirmizi] (-7/3,0) circle (2.5pt) node[above] {\small \(-\frac{7}{3}\)};
\fill[cozumyesil] (11/4,0) circle (2.5pt) node[above] {\small \(\frac{11}{4}\)};
\end{tikzpicture}
\end{center}
\(0\)'a uzaklıklar: \(\frac{7}{3}=\frac{28}{12}<\frac{33}{12}=\frac{11}{4}\) \(\Rightarrow\) \(-\frac{7}{3}\) daha yakın.}

\subsection{Toplama ve çıkarma}
\begin{frame}{Kesirlerde Toplama ve Çıkarma}
\begin{block}{Kural}
\begin{itemize}
\item Paydalar eşitse: \(\dfrac{a}{c}\pm\dfrac{b}{c}=\dfrac{a\pm b}{c}\)
\item Paydalar farklıysa önce paydalar eşitlenir. En küçük ortak payda \(=\) paydaların \textbf{EKOK}'u.
\item Tam sayılı kesirler önce bileşik kesre çevrilir.
\item Sonuç en sade biçimde yazılır.
\end{itemize}
\end{block}
\begin{alertblock}{Sık yapılan hata}
\(\dfrac{1}{2}+\dfrac{1}{3}\neq\dfrac{1+1}{2+3}=\dfrac{2}{5}\). \quad Doğrusu: \(\dfrac{3}{6}+\dfrac{2}{6}=\dfrac{5}{6}\). \quad Paydalar toplanmaz!
\end{alertblock}
\end{frame}

\soru{Toplama -- Çıkarma: Örnek 6}
{\(\dfrac{2}{3}+\dfrac{3}{4}-\dfrac{5}{6}\) işleminin sonucu kaçtır?}
{\(\text{EKOK}(3,4,6)=12\):
\[\frac{2}{3}+\frac{3}{4}-\frac{5}{6}=\frac{8}{12}+\frac{9}{12}-\frac{10}{12}=\frac{7}{12}\]}

\soru{Toplama -- Çıkarma: Örnek 7}
{\(\dfrac{5}{12}+\dfrac{7}{18}\) işleminin sonucu kaçtır?}
{\(12=2^2\cdot 3\), \ \(18=2\cdot 3^2\) \(\Rightarrow\) \(\text{EKOK}=2^2\cdot 3^2=36\):
\[\frac{5}{12}+\frac{7}{18}=\frac{15}{36}+\frac{14}{36}=\frac{29}{36}\]
(\(12\cdot 18=216\) de ortak paydadır ama sayılar büyür.)}

\soru{Tam Sayılı Kesirler -- Örnek 8}
{\(3\dfrac{1}{4}-1\dfrac{2}{3}\) işleminin sonucu kaçtır?}
{\[3\frac{1}{4}-1\frac{2}{3}=\frac{13}{4}-\frac{5}{3}=\frac{39}{12}-\frac{20}{12}=\frac{19}{12}=1\frac{7}{12}\]}

\subsection{Çarpma, bölme, merdiven kesirler}
\begin{frame}{Kesirlerde Çarpma, Bölme ve Merdiven Kesirler}
\begin{block}{Kural}
\[\frac{a}{b}\cdot\frac{c}{d}=\frac{a\cdot c}{b\cdot d},\qquad \frac{a}{b}\div\frac{c}{d}=\frac{a}{b}\cdot\frac{d}{c}\qquad(b,c,d\neq 0)\]
\vspace{-6pt}
\begin{itemize}
\item Çarpmadan önce çapraz sadeleştirme yapmak işlemi kısaltır.
\item Bölmede: birinci kesir aynen yazılır, ikinci kesir \textbf{ters çevrilip} çarpılır.
\item \textbf{Merdiven kesir}: En alttaki (en içteki) kesirden başlanarak yukarı doğru hesaplanır.
\end{itemize}
\end{block}
\begin{alertblock}{Sık yapılan hata: ana kesir çizgisi}
\[\frac{3}{\ \dfrac{4}{5}\ }=3\cdot\frac{5}{4}=\frac{15}{4}\qquad\text{ama}\qquad \frac{\ \dfrac{3}{4}\ }{5}=\frac{3}{4}\cdot\frac{1}{5}=\frac{3}{20}\]
\end{alertblock}
\end{frame}

\soru{Çarpma -- Örnek 9}
{\(\dfrac{14}{15}\cdot\dfrac{25}{21}\) işleminin sonucu kaçtır?}
{Çapraz sadeleştirme: \(14\) ile \(21\) (\(\div 7\)), \ \(25\) ile \(15\) (\(\div 5\)):
\[\frac{14}{15}\cdot\frac{25}{21}=\frac{2}{3}\cdot\frac{5}{3}=\frac{10}{9}\]}

\soru{Bölme -- Örnek 10}
{\(\dfrac{4}{9}\div\dfrac{8}{15}\) \ ve \ \(\left(-\dfrac{3}{5}\right)\div 6\) işlemlerinin sonuçları nedir?}
{\[\frac{4}{9}\div\frac{8}{15}=\frac{4}{9}\cdot\frac{15}{8}=\frac{1}{3}\cdot\frac{5}{2}=\frac{5}{6}\]
\[\left(-\frac{3}{5}\right)\div 6=-\frac{3}{5}\cdot\frac{1}{6}=-\frac{1}{10}\]}

\soru{Merdiven Kesir -- Örnek 11}
{\(\dfrac{\dfrac{2}{3}-\dfrac{1}{2}}{\dfrac{1}{4}+\dfrac{1}{6}}\) işleminin sonucu kaçtır?}
{Pay: \(\frac{4}{6}-\frac{3}{6}=\frac{1}{6}\); \quad Payda: \(\frac{3}{12}+\frac{2}{12}=\frac{5}{12}\)
\[\frac{1}{6}\div\frac{5}{12}=\frac{1}{6}\cdot\frac{12}{5}=\frac{2}{5}\]}

\soru{Merdiven Kesir -- Örnek 12}
{\(2-\cfrac{1}{2-\cfrac{1}{2-\cfrac{1}{2}}}\) işleminin sonucu kaçtır?}
{En alttan başla:
\[2-\tfrac{1}{2}=\tfrac{3}{2}\ \rightarrow\ \tfrac{1}{3/2}=\tfrac{2}{3}\ \rightarrow\ 2-\tfrac{2}{3}=\tfrac{4}{3}\ \rightarrow\ \tfrac{1}{4/3}=\tfrac{3}{4}\ \rightarrow\ 2-\tfrac{3}{4}=\tfrac{5}{4}\]}

\begin{frame}{Ders 1 -- Sınıfta Siz Çözün}
\begin{columns}[t]
\column{0.58\textwidth}
\begin{enumerate}
\item \(4\frac{2}{7}\) sayısını bileşik kesir yazınız.
\item \(\dfrac{90}{150}\) kesrinin en sade biçimi? \(\star\)
\item \(\dfrac{5}{6}-\dfrac{3}{8}=?\) \(\star\)
\item \(\dfrac{9}{10}\div\dfrac{3}{5}=?\)
\item \(1-\cfrac{1}{1-\cfrac{1}{3}}=?\) \(\star\)
\end{enumerate}
\column{0.36\textwidth}
\pause
\textbf{Cevaplar:}
\begin{enumerate}
\item \(\frac{30}{7}\)
\item \(\frac{3}{5}\)
\item \(\frac{11}{24}\)
\item \(\frac{3}{2}\)
\item \(-\frac{1}{2}\)
\end{enumerate}
\end{columns}
\vspace{4pt}{\footnotesize \(\star\) işaretli soruların adım adım çözümü sonraki slaytlardadır.}
\end{frame}

% ---- Ders 1 alıştırma çözümleri ----
\soru{Ders 1 -- Alıştırma Çözümü (Soru 2)}
{\(\dfrac{90}{150}\) kesrinin en sade biçimi nedir?}
{Pay ve payda EBOB'larına bölünür:\\[3pt]
\(90=2\cdot 3^2\cdot 5\), \ \(150=2\cdot 3\cdot 5^2\) \(\Rightarrow\) \(\text{EBOB}(90,150)=2\cdot 3\cdot 5=30\)
\[\frac{90}{150}=\frac{90\div 30}{150\div 30}=\frac{3}{5}\]}

\soru{Ders 1 -- Alıştırma Çözümü (Soru 3)}
{\(\dfrac{5}{6}-\dfrac{3}{8}\) işleminin sonucu kaçtır?}
{Paydalar eşitlenir: \(\text{EKOK}(6,8)=24\).
\[\frac{5}{6}=\frac{20}{24},\qquad \frac{3}{8}=\frac{9}{24}\]
\[\frac{20}{24}-\frac{9}{24}=\frac{11}{24}\qquad(\text{\(11\) ile \(24\) aralarında asal: en sade})\]}

\soru{Ders 1 -- Alıştırma Çözümü (Soru 5)}
{\(1-\cfrac{1}{1-\cfrac{1}{3}}\) işleminin sonucu kaçtır?}
{En alttan başla:
\[1-\tfrac{1}{3}=\tfrac{2}{3}\ \rightarrow\ \frac{1}{2/3}=\frac{3}{2}\ \rightarrow\ 1-\frac{3}{2}=-\frac{1}{2}\]
Sonuç negatif olabilir: \(-\dfrac{1}{2}\).}

% =====================================================
\section{Ders 2: Ondalık Sayılar}
\begin{frame}
\centering\Large\color{anamavi}\textbf{Ders 2}\\[6pt]\normalsize Ondalık Sayılar, Devirli Ondalıklar ve Bilimsel Gösterim
\end{frame}

\subsection{Ondalık ve kesir}
\begin{frame}{Ondalık Sayılar}
\begin{block}{Basamaklar ve dönüşüm}
\(4{,}375\): \ \(4\) birler, \(3\) onda birler, \(7\) yüzde birler, \(5\) binde birler basamağında.
\begin{itemize}
\item \textbf{Ondalık \(\to\) kesir:} virgülsüz sayı / (\(1\) ve virgülden sonraki basamak sayısı kadar \(0\)):
\ \(0{,}25=\dfrac{25}{100}=\dfrac{1}{4}\), \quad \(4{,}375=\dfrac{4375}{1000}=\dfrac{35}{8}\)
\item \textbf{Kesir \(\to\) ondalık:} paydayı \(10,100,1000,\dots\) yap ya da payı paydaya böl: \ \(\dfrac{3}{4}=\dfrac{75}{100}=0{,}75\)
\end{itemize}
\end{block}
\begin{alertblock}{Ne zaman sonlu ondalık?}
En sade kesrin paydasının asal çarpanları yalnızca \(2\) ve/veya \(5\) ise ondalık açılım sonludur:
\(\frac{3}{40}=0{,}075\). \ Aksi halde devirlidir: \(\frac{1}{3}=0{,}333\ldots\), \ \(\frac{5}{12}=0{,}41666\ldots\)
\end{alertblock}
\end{frame}

\soru{Ondalık \(\to\) Kesir: Örnek 13}
{\(0{,}25\), \ \(1{,}6\) \ ve \ \(0{,}045\) sayılarını en sade kesir olarak yazınız.}
{\[0{,}25=\frac{25}{100}=\frac{1}{4},\qquad 1{,}6=\frac{16}{10}=\frac{8}{5}\]
\[0{,}045=\frac{45}{1000}=\frac{9}{200}\]}

\soru{Kesir \(\to\) Ondalık: Örnek 14}
{\(\dfrac{7}{8}\) \ ve \ \(\dfrac{3}{20}\) kesirlerini ondalık sayı olarak yazınız.}
{\(8=2^3\): \ \(5^3=125\) ile genişlet:
\[\frac{7}{8}=\frac{7\cdot 125}{8\cdot 125}=\frac{875}{1000}=0{,}875\]
\[\frac{3}{20}=\frac{3\cdot 5}{20\cdot 5}=\frac{15}{100}=0{,}15\]}

\subsection{Dört işlem}
\begin{frame}{Ondalık Sayılarda Dört İşlem}
\begin{block}{Kural}
\begin{itemize}
\item \textbf{Toplama/çıkarma:} Virgüller alt alta gelecek şekilde yazılır (eksik basamağa \(0\)).
\item \textbf{Çarpma:} Virgül yokmuş gibi çarpılır; sonuçta virgülden sonra, çarpanlardaki ondalık basamak sayılarının \textbf{toplamı} kadar basamak ayrılır.
\item \textbf{Bölme:} Pay ve payda aynı \(10^k\) ile genişletilerek virgüllerden kurtulunur:\ \(\dfrac{0{,}36}{0{,}012}=\dfrac{360}{12}\)
\item \(10,100,1000\) ile çarpma: virgül sağa \(1,2,3\) basamak kayar; bölme: sola kayar.
\end{itemize}
\end{block}
\begin{alertblock}{Sık yapılan hata}
\(0{,}3\cdot 0{,}2=0{,}06\) \ (\(0{,}6\) değil); \quad \(0{,}2^{\,2}=0{,}04\) \ (\(0{,}4\) değil).
\end{alertblock}
\end{frame}

\soru{Toplama -- Çıkarma: Örnek 15}
{\(3{,}75+12{,}4-0{,}08\) işleminin sonucu kaçtır?}
{Virgülleri hizala:
\[3{,}75+12{,}40=16{,}15\]
\[16{,}15-0{,}08=16{,}07\]}

\soru{Çarpma -- Örnek 16}
{\(2{,}4\cdot 0{,}35\), \ \(0{,}06\cdot 0{,}5\), \ \(3{,}47\cdot 100\), \ \(52{,}1\div 1000\) işlemlerinin sonuçları nedir?}
{\(24\cdot 35=840\), ondalık basamak \(1+2=3\): \ \(2{,}4\cdot 0{,}35=0{,}840=0{,}84\)\\[3pt]
\(6\cdot 5=30\), ondalık basamak \(2+1=3\): \ \(0{,}06\cdot 0{,}5=0{,}030=0{,}03\)\\[3pt]
\(3{,}47\cdot 100=347\) \ (virgül \(2\) sağa), \quad \(52{,}1\div 1000=0{,}0521\) \ (virgül \(3\) sola)}

\soru{Bölme -- Örnek 17}
{\(\dfrac{0{,}36}{0{,}012}\) \ ve \ \(\dfrac{0{,}24}{0{,}6}+\dfrac{1{,}5}{0{,}05}\) işlemlerinin sonuçları nedir?}
{Pay ve paydayı \(1000\) ile genişlet: \quad \(\dfrac{0{,}36}{0{,}012}=\dfrac{360}{12}=30\)\\[4pt]
\(\dfrac{0{,}24}{0{,}6}=\dfrac{24}{60}=0{,}4\), \quad \(\dfrac{1{,}5}{0{,}05}=\dfrac{150}{5}=30\) \quad \(\Rightarrow\) \quad \(0{,}4+30=30{,}4\)}

\subsection{Devirli ondalık sayılar}
\begin{frame}{Devirli Ondalık Sayılar}
\begin{block}{Kesre çevirme kuralı}
\[\text{Kesir}=\frac{\text{(virgülsüz tüm sayı)}-\text{(devretmeyen kısım)}}{9\ldots9\,0\ldots0}\]
\vspace{-4pt}
Paydada: \textbf{devreden} basamak sayısı kadar \(9\), ardından \textbf{virgülden sonra devretmeyen} basamak sayısı kadar \(0\).
\[0{,}\ab{3}=\frac{3-0}{9}=\frac{1}{3},\qquad 1{,}2\ab{3}=\frac{123-12}{90}=\frac{111}{90}=\frac{37}{30}\]
\end{block}
\begin{alertblock}{Dikkat}
\(0{,}\ab{9}=\dfrac{9-0}{9}=1\): \ \(0{,}999\ldots\) ile \(1\) \textbf{aynı sayıdır}. \quad Devirli sayılarla işlem yapmadan önce kesre çevirin.
\end{alertblock}
\end{frame}

\soru{Devirli Ondalık -- Örnek 18}
{\(0{,}\ab{3}\) \ ve \ \(0{,}\ab{27}\) sayılarını en sade kesir olarak yazınız.}
{\[0{,}\ab{3}=\frac{3}{9}=\frac{1}{3}\]
İki basamak devrediyor \(\Rightarrow\) paydada iki tane \(9\):
\[0{,}\ab{27}=\frac{27}{99}=\frac{3}{11}\]}

\soru{Devirli Ondalık -- Örnek 19}
{\(1{,}2\ab{3}\) sayısını en sade kesir olarak yazınız.}
{Tüm sayı \(123\), devretmeyen kısım \(12\).\\
Devreden: \(1\) basamak (\(3\)) \(\Rightarrow\) bir \(9\); \ virgülden sonra devretmeyen: \(1\) basamak (\(2\)) \(\Rightarrow\) bir \(0\).
\[1{,}2\ab{3}=\frac{123-12}{90}=\frac{111}{90}=\frac{37}{30}\]}

\soru{Devirli Ondalık -- Örnek 20}
{\(0{,}1\ab{6}\) \ ve \ \(2{,}\ab{45}\) sayılarını en sade kesir olarak yazınız.}
{\[0{,}1\ab{6}=\frac{16-1}{90}=\frac{15}{90}=\frac{1}{6}\]
\[2{,}\ab{45}=\frac{245-2}{99}=\frac{243}{99}=\frac{27}{11}\]}

\soru{Devirli Ondalık -- Örnek 21}
{\(2{,}4\ab{9}\) sayısı hangi sonlu ondalık sayıya eşittir?}
{\[2{,}4\ab{9}=\frac{249-24}{90}=\frac{225}{90}=\frac{5}{2}=2{,}5\]
Genel olarak: sondaki \(\ab{9}\) bir önceki basamağı \(1\) artırır (\(0{,}\ab{9}=1\) gibi).}

\soru{Devirli Ondalık -- Örnek 22}
{\(\dfrac{0{,}\ab{8}}{0{,}\ab{4}}+0{,}\ab{6}\cdot 1{,}5\) işleminin sonucu kaçtır?}
{\(0{,}\ab{8}=\frac{8}{9}\), \ \(0{,}\ab{4}=\frac{4}{9}\), \ \(0{,}\ab{6}=\frac{2}{3}\), \ \(1{,}5=\frac{3}{2}\)
\[\frac{8}{9}\div\frac{4}{9}+\frac{2}{3}\cdot\frac{3}{2}=2+1=3\]}

\subsection{Bilimsel gösterim}
\begin{frame}{Bilimsel Gösterim (Giriş)}
\begin{block}{Tanım}
\[a\cdot 10^{n},\qquad 1\le |a|<10,\quad n\in\mathbb{Z}\]
\(10^{1}=10,\ 10^{2}=100,\ 10^{3}=1000,\ \dots\) \qquad \(10^{0}=1\) \qquad \(10^{-1}=0{,}1,\ 10^{-2}=0{,}01,\ 10^{-3}=0{,}001,\ \dots\)
\begin{itemize}
\item Büyük sayı: virgül \(n\) basamak sola kaydırılırsa \(\cdot\,10^{n}\) yazılır.
\item Küçük sayı: virgül \(n\) basamak sağa kaydırılırsa \(\cdot\,10^{-n}\) yazılır.
\end{itemize}
\end{block}
\begin{alertblock}{Dikkat}
\(34\cdot 10^{5}\) bilimsel gösterim değildir (\(34\ge 10\)); doğrusu \(3{,}4\cdot 10^{6}\). \ Üslü sayıların özellikleri 9. haftada.
\end{alertblock}
\end{frame}

\soru{Bilimsel Gösterim -- Örnek 23}
{\(3\,400\,000\) \ ve \ \(0{,}00052\) sayılarını bilimsel gösterimle yazınız.}
{\(3\,400\,000\): virgül \(6\) basamak sola \(\Rightarrow\) \(3{,}4\cdot 10^{6}\)\\[4pt]
\(0{,}00052\): virgül \(4\) basamak sağa \(\Rightarrow\) \(5{,}2\cdot 10^{-4}\)\\[4pt]
Kontrol: \(5{,}2\cdot 0{,}0001=0{,}00052\).}

\begin{frame}{Ders 2 -- Sınıfta Siz Çözün}
\begin{columns}[t]
\column{0.58\textwidth}
\begin{enumerate}
\item \(0{,}125\) sayısını en sade kesir yazınız.
\item \(4{,}2\cdot 0{,}05=?\)
\item \(\dfrac{0{,}72}{0{,}018}=?\) \(\star\)
\item \(0{,}2\ab{7}\) sayısını en sade kesir yazınız. \(\star\)
\item \(0{,}0045\cdot 10^{5}\) sayısını bilimsel gösterimle yazınız. \(\star\)
\end{enumerate}
\column{0.36\textwidth}
\pause
\textbf{Cevaplar:}
\begin{enumerate}
\item \(\frac{1}{8}\)
\item \(0{,}21\)
\item \(40\)
\item \(\frac{27-2}{90}=\frac{5}{18}\)
\item \(450=4{,}5\cdot 10^{2}\)
\end{enumerate}
\end{columns}
\vspace{4pt}{\footnotesize \(\star\) işaretli soruların adım adım çözümü sonraki slaytlardadır.}
\end{frame}

% ---- Ders 2 alıştırma çözümleri ----
\soru{Ders 2 -- Alıştırma Çözümü (Soru 3)}
{\(\dfrac{0{,}72}{0{,}018}\) işleminin sonucu kaçtır?}
{En çok ondalık basamak: \(0{,}018\) (\(3\) basamak) \(\Rightarrow\) pay ve payda \(1000\) ile genişletilir:
\[\frac{0{,}72}{0{,}018}=\frac{720}{18}=40\]
Kontrol: \(0{,}018\cdot 40=0{,}72\) \checkmark}

\soru{Ders 2 -- Alıştırma Çözümü (Soru 4)}
{\(0{,}2\ab{7}\) sayısını en sade kesir olarak yazınız.}
{Virgülsüz tüm sayı: \(27\); devretmeyen kısım: \(2\).\\[3pt]
Payda: bir basamak devrediyor \(\to\) bir \(9\); virgülden sonra bir basamak devretmiyor \(\to\) bir \(0\): \(90\).
\[0{,}2\ab{7}=\frac{27-2}{90}=\frac{25}{90}=\frac{5}{18}\]}

\soru{Ders 2 -- Alıştırma Çözümü (Soru 5)}
{\(0{,}0045\cdot 10^{5}\) sayısını bilimsel gösterimle yazınız.}
{\(10^{5}\) ile çarpmak virgülü \(5\) basamak sağa kaydırır: \(0{,}0045\cdot 10^{5}=450\).\\[3pt]
\(450\): virgül \(2\) basamak sola \(\Rightarrow\) \(4{,}5\cdot 10^{2}\) \quad (\(1\le 4{,}5<10\) \checkmark)}

% =====================================================
\section{Ders 3: Sıralama, Karma İşlem, Kesir Problemleri}
\begin{frame}
\centering\Large\color{anamavi}\textbf{Ders 3}\\[6pt]\normalsize Rasyonel Sayılarda Sıralama, Karma İşlemler ve Kesir Problemleri
\end{frame}

\subsection{Sıralama}
\begin{frame}{Rasyonel Sayılarda Sıralama}
\begin{block}{Yöntemler (\(a,b,c,d>0\))}
\begin{itemize}
\item \textbf{Paydalar eşit:} payı büyük olan büyüktür. \quad \textbf{Paylar eşit:} paydası küçük olan büyüktür.
\item \textbf{Payda (veya pay) eşitleme:} EKOK ile ortak payda / ortak pay bulunur.
\item \textbf{Ondalığa çevirme:} özellikle ondalık ve devirli sayılar karışıksa.
\item \textbf{\(1\)'e fark:} basit kesirlerde (\(\frac{13}{14}=1-\frac{1}{14}\)) \(1\)'den eksiği az olan büyüktür; bileşik kesirlerde \(1\)'den fazlası çok olan büyüktür.
\item \textbf{Çapraz çarpım:} \(\dfrac{a}{b}<\dfrac{c}{d}\iff a\cdot d<b\cdot c\).
\end{itemize}
\end{block}
\begin{alertblock}{Sık yapılan hatalar}
Negatiflerde sıra \textbf{ters döner}: \(\frac{1}{3}<\frac{1}{2}\) ama \(-\frac{1}{2}<-\frac{1}{3}\). \quad Çapraz çarpımda paydalar pozitif olmalı.
\end{alertblock}
\end{frame}

\soru{Sıralama -- Örnek 24}
{\(a=\dfrac{5}{9},\quad b=\dfrac{5}{7},\quad c=\dfrac{5}{12}\) \ sayılarını küçükten büyüğe sıralayınız.}
{Paylar eşit (\(5\)) \(\Rightarrow\) paydası küçük olan büyük.\\[3pt]
Paydalar: \(7<9<12\) \(\Rightarrow\) \(b>a>c\)
\[c<a<b\]}

\soru{Sıralama -- Örnek 25}
{\(\dfrac{2}{3},\quad \dfrac{3}{4},\quad \dfrac{5}{8}\) \ sayılarını küçükten büyüğe sıralayınız.}
{\(\text{EKOK}(3,4,8)=24\):
\[\frac{2}{3}=\frac{16}{24},\qquad \frac{3}{4}=\frac{18}{24},\qquad \frac{5}{8}=\frac{15}{24}\]
\[\frac{5}{8}<\frac{2}{3}<\frac{3}{4}\]}

\soru{Sıralama -- Örnek 26}
{\(a=0{,}65,\quad b=\dfrac{2}{3},\quad c=0{,}6\ab{5}\) \ sayılarını küçükten büyüğe sıralayınız.}
{Ondalığa çevir, basamak basamak karşılaştır:
\[a=0{,}6500\ldots,\qquad c=0{,}6555\ldots,\qquad b=0{,}6666\ldots\]
\[a<c<b\]}

\soru{Sıralama -- Örnek 27}
{\(\dfrac{13}{14},\quad \dfrac{15}{16},\quad \dfrac{17}{18}\) \ sayılarını küçükten büyüğe sıralayınız.}
{\(1\)'e olan farklar: \ \(\dfrac{1}{14},\ \dfrac{1}{16},\ \dfrac{1}{18}\). \quad \(\dfrac{1}{18}<\dfrac{1}{16}<\dfrac{1}{14}\)\\[3pt]
\(1\)'e en yakın olan en büyüktür:
\[\frac{13}{14}<\frac{15}{16}<\frac{17}{18}\]}

\soru{Sıralama -- Örnek 28}
{\(\dfrac{17}{23}\) \ ile \ \(\dfrac{13}{18}\) \ sayılarından hangisi büyüktür?}
{Çapraz çarpım (paydalar pozitif):
\[17\cdot 18=306,\qquad 13\cdot 23=299\]
\(306>299\) \(\Rightarrow\) \(\dfrac{17}{23}>\dfrac{13}{18}\)}

\soru{Negatif Rasyoneller -- Örnek 29}
{\(-\dfrac{2}{3},\quad -\dfrac{3}{4},\quad -\dfrac{5}{8}\) \ sayılarını küçükten büyüğe sıralayınız.}
{Pozitifleri (Örnek 25): \ \(\frac{5}{8}<\frac{2}{3}<\frac{3}{4}\).\\[3pt]
Negatiflerde sıra ters döner:
\[-\frac{3}{4}<-\frac{2}{3}<-\frac{5}{8}\]}

\soru{Arasında Sayı -- Örnek 30}
{\(\dfrac{1}{3}\) ile \(\dfrac{1}{2}\) arasında bir rasyonel sayı yazınız. \(\dfrac{3}{7}\) ile \(\dfrac{4}{7}\) arasında olup paydası \(14\) olan kesir hangisidir?}
{Ortalama arada kalır: \(\dfrac{\frac{1}{3}+\frac{1}{2}}{2}=\dfrac{5/6}{2}=\dfrac{5}{12}\) \quad (\(\frac{4}{12}<\frac{5}{12}<\frac{6}{12}\))\\[4pt]
\(\dfrac{3}{7}=\dfrac{6}{14}\), \ \(\dfrac{4}{7}=\dfrac{8}{14}\) \(\Rightarrow\) arada: \(\dfrac{7}{14}=\dfrac{1}{2}\)\\[3pt]
{\normalsize İki farklı rasyonel sayı arasında sonsuz tane rasyonel sayı vardır.}}

\subsection{Karma işlem ve problemler}
\begin{frame}{Karma İşlemler ve Kesir Problemleri}
\begin{block}{Kural}
\begin{itemize}
\item İşlem önceliği (1. hafta) aynen geçerlidir: parantez \(\to\) \(\times,\div\) \(\to\) \(+,-\), soldan sağa.
\item Ondalık ve devirli sayılar önce kesre çevrilir.
\item Bir bütünün \(\dfrac{a}{b}\)'si \(=\) bütün \(\cdot\dfrac{a}{b}\) \quad (ör. \(60\)'ın \(\frac{3}{4}\)'ü \(=45\))
\end{itemize}
\end{block}
\begin{alertblock}{Sık yapılan hata: ``tamamının'' mı, ``kalanın'' mı?}
``Tamamının \(\frac{1}{3}\)'ü'' başlangıçtaki bütüne göre; ``kalanın \(\frac{1}{3}\)'ü'' o anda kalan miktara göre hesaplanır.
\end{alertblock}
\end{frame}

\soru{Karma İşlem -- Örnek 31}
{\(\left(1-\dfrac{1}{2}\right)\left(1-\dfrac{1}{3}\right)\left(1-\dfrac{1}{4}\right)\cdots\left(1-\dfrac{1}{10}\right)\) çarpımı kaçtır?}
{Her parantezi hesapla:
\[\frac{1}{2}\cdot\frac{2}{3}\cdot\frac{3}{4}\cdots\frac{9}{10}\]
Her pay bir önceki paydayı sadeleştirir; geriye ilk pay ve son payda kalır:
\[=\frac{1}{10}\]}

\soru{Kesir Problemi -- Örnek 32}
{\(360\) TL'nin \(\dfrac{2}{3}\)'ü ile \(\dfrac{5}{8}\)'i arasındaki fark kaç TL'dir?}
{\[360\cdot\frac{2}{3}=240,\qquad 360\cdot\frac{5}{8}=225\]
\[240-225=15\ \text{TL}\]
Kısa yol: \(360\cdot\left(\frac{2}{3}-\frac{5}{8}\right)=360\cdot\frac{1}{24}=15\)}

\soru{Kesir Problemi -- Örnek 33}
{Bir deponun \(\frac{3}{4}\)'ü doludur. (a) Deponun \textbf{tamamının} \(\frac{1}{3}\)'ü kadar su kullanılırsa, (b) içindeki suyun \(\frac{1}{3}\)'ü kullanılırsa depo kaçta kaçı dolu kalır?}
{(a) \(\dfrac{3}{4}-\dfrac{1}{3}=\dfrac{9}{12}-\dfrac{4}{12}=\dfrac{5}{12}\)\\[4pt]
(b) Suyun \(\frac{2}{3}\)'ü kalır: \quad \(\dfrac{3}{4}\cdot\dfrac{2}{3}=\dfrac{1}{2}\)\\[4pt]
Aynı ``\(\frac{1}{3}\)'', farklı bütün \(\Rightarrow\) farklı sonuç.}

\soru{Kesir Problemi -- Örnek 34}
{\(120\) sayfalık bir kitabın ilk gün \(\frac{1}{4}\)'ü, ikinci gün kalan sayfaların \(\frac{2}{5}\)'i okunuyor. Okunmamış kaç sayfa kalır?}
{1. gün: \(120\cdot\frac{1}{4}=30\) sayfa \(\Rightarrow\) kalan \(90\)\\[3pt]
2. gün: \(90\cdot\frac{2}{5}=36\) sayfa \(\Rightarrow\) kalan \(90-36=54\)\\[3pt]
Kısa yol: \(120\cdot\frac{3}{4}\cdot\frac{3}{5}=54\) sayfa.}

\begin{frame}{Ders 3 -- Sınıfta Siz Çözün}
\begin{columns}[t]
\column{0.6\textwidth}
\begin{enumerate}
\item \(\frac{4}{5},\ \frac{5}{6},\ \frac{6}{7}\) küçükten büyüğe?
\item \(-\frac{1}{2},\ -\frac{2}{5},\ -\frac{3}{7}\) küçükten büyüğe? \(\star\)
\item \(\frac{2}{9}\) ile \(\frac{1}{3}\)'ün ortalaması? \(\star\)
\item \(240\) öğrencinin \(\frac{3}{8}\)'i kız ise kaç erkek öğrenci vardır?
\item \(\left(1+\frac{1}{2}\right)\left(1+\frac{1}{3}\right)\cdots\left(1+\frac{1}{9}\right)=?\) \(\star\)
\end{enumerate}
\column{0.36\textwidth}
\pause
\textbf{Cevaplar:}
\begin{enumerate}
\item \(\frac{4}{5}<\frac{5}{6}<\frac{6}{7}\)
\item \(-\frac{1}{2}<-\frac{3}{7}<-\frac{2}{5}\)
\item \(\frac{5}{18}\)
\item \(150\)
\item \(\frac{10}{2}=5\)
\end{enumerate}
\end{columns}
\vspace{4pt}{\footnotesize \(\star\) işaretli soruların adım adım çözümü sonraki slaytlardadır.}
\end{frame}

% ---- Ders 3 alıştırma çözümleri ----
\soru{Ders 3 -- Alıştırma Çözümü (Soru 2)}
{\(-\dfrac{1}{2},\ -\dfrac{2}{5},\ -\dfrac{3}{7}\) sayılarını küçükten büyüğe sıralayınız.}
{Önce pozitiflerini sıralayalım: \(\text{EKOK}(2,5,7)=70\).
\[\frac{1}{2}=\frac{35}{70},\quad \frac{2}{5}=\frac{28}{70},\quad \frac{3}{7}=\frac{30}{70}\ \Rightarrow\ \frac{2}{5}<\frac{3}{7}<\frac{1}{2}\]
Negatiflerde sıra ters döner: \ \(-\dfrac{1}{2}<-\dfrac{3}{7}<-\dfrac{2}{5}\)}

\soru{Ders 3 -- Alıştırma Çözümü (Soru 3)}
{\(\dfrac{2}{9}\) ile \(\dfrac{1}{3}\) sayılarının aritmetik ortalaması kaçtır?}
{\(\dfrac{1}{3}=\dfrac{3}{9}\) \(\Rightarrow\) toplam: \(\dfrac{2}{9}+\dfrac{3}{9}=\dfrac{5}{9}\)
\[\text{Ortalama}=\frac{5}{9}\div 2=\frac{5}{9}\cdot\frac{1}{2}=\frac{5}{18}\]
Kontrol: \(\dfrac{4}{18}<\dfrac{5}{18}<\dfrac{6}{18}\) \checkmark}

\soru{Ders 3 -- Alıştırma Çözümü (Soru 5)}
{\(\left(1+\frac{1}{2}\right)\left(1+\frac{1}{3}\right)\cdots\left(1+\frac{1}{9}\right)\) çarpımı kaçtır?}
{Her çarpan: \(1+\dfrac{1}{n}=\dfrac{n+1}{n}\).
\[\frac{3}{2}\cdot\frac{4}{3}\cdot\frac{5}{4}\cdots\frac{10}{9}\]
Her pay bir sonraki paydayla sadeleşir; yalnız ilk payda ve son pay kalır:
\[\frac{10}{2}=5\]}

% =====================================================
\begin{frame}{Özet}
\small
\begin{columns}[t]
\column{0.32\textwidth}
\begin{block}{Ders 1}
\begin{itemize}
\item \(\frac{a}{b}\), \(b\neq 0\); \ \(\frac{a}{0}\) tanımsız
\item Sadeleştir: EBOB'a böl
\item \(\pm\): EKOK ile payda eşitle
\item \(\div\): ters çevir, çarp
\item Merdiven: en alttan başla
\end{itemize}
\end{block}
\column{0.32\textwidth}
\begin{block}{Ders 2}
\begin{itemize}
\item \(\times\): ondalık basamaklar toplanır
\item \(\div\): \(10^{k}\) ile genişlet
\item Devirli: \(\dfrac{\text{tüm}-\text{devretmeyen}}{9\ldots 90\ldots 0}\)
\item \(0{,}\ab{9}=1\)
\item \(a\cdot 10^{n}\), \(1\le|a|<10\)
\end{itemize}
\end{block}
\column{0.32\textwidth}
\begin{block}{Ders 3}
\begin{itemize}
\item Payda/pay eşitle, ondalığa çevir
\item \(1\)'e fark, çapraz çarpım
\item Negatiflerde sıra ters
\item Ortalama arada kalır
\item ``Tamamının'' \(\neq\) ``kalanın''
\end{itemize}
\end{block}
\end{columns}
\vspace{2pt}
\begin{alertblock}{Gelecek hafta}
Birinci dereceden bir bilinmeyenli denklemler (kesirli denklemlerde bu haftanın payda eşitleme becerisi gerekecek)
\end{alertblock}
\end{frame}

\begin{frame}
\centering\Large İyi çalışmalar
\end{frame}

\end{document}
